\section{Cross-Domain Case Walkthroughs}

Table~\ref{tab:case-walkthroughs} turns the three running examples into disturbance tests. In each row, start from the same arrivals and state, vary the named decision, and follow both the immediate beneficiaries and the work exposed to its cost.

\begin{table}[H]
\small
\centering
\setlength{\tabcolsep}{4pt}
\caption{Events and measurements for the three running examples.}
\label{tab:case-walkthroughs}
\begin{tabular}{@{}p{0.16\textwidth}p{0.38\textwidth}p{0.40\textwidth}@{}}
\toprule
Case & Event sequence and comparison & Observable consequences \\
\midrule
Serving &
A leaves a valid, unpinned KV prefix. When B needs pages, reclaim the prefix or make B wait. Revisit A with the original input, model, and adapter~\cite{kwon2023pagedattention,hu2024memserve,liu2025lmcache}. &
Record pages freed and B's first-token delay. Record A's first-token delay and whether its prefix was loaded or rebuilt. \\
Shard handoff &
Move independently assignable keys. Inject failure after copying but before the ownership commit, then repeat just after the commit~\cite{carbone2015flink,nasir2019megaphone,mao2025spacker}. &
Check the authoritative epoch and last committed input position. Track duplicate or missed updates, write pause, both workers' queues, and recovery time. \\
Index update &
Delete a document while queries arrive. Compare repair schedules under the same update trace and enforce each query's minimum visible version~\cite{singh2021freshdiskann,xu2023spfresh,wang2026candor}. &
Track visibility lag, recall, and latency during maintenance. Report how often a query waits or is redirected because its local replica has an older version. \\
\bottomrule
\end{tabular}
\end{table}
